% TOP_PROBLEM: {"id": 9700019, "title": "Largest common subcladogram exponents", "category_id": 19, "category": {"id": 19, "name": "probability", "display_name": "Probability", "description": "Problems involving probability theory, stochastic processes, and random structures.", "slug": "probability", "order_index": 19, "created_at": "2026-07-31T15:26:25.671Z"}, "status": "open", "problem_number": "AMR-096-0019", "set_id": 13}
% FIRST_POSTED: 2026-10-01
% ENTRY_KIND: statement_only
% UnsolvedMath ID 9700019; source catalogue code AMR-096-0019
% !TeX program = lualatex
% Definition and sources only; this file is not an LLM solution attempt.
\documentclass[11pt,a4paper]{article}
\usepackage[margin=25mm]{geometry}
\usepackage{fontspec}
\setmainfont{lmroman10-regular.otf}[BoldFont=lmroman10-bold.otf,
 ItalicFont=lmroman10-italic.otf,BoldItalicFont=lmroman10-bolditalic.otf]
\usepackage{amsmath,amssymb,mathtools}
\usepackage{unicode-math}
\setmathfont{latinmodern-math.otf}
\AtBeginDocument{\let\setminus\smallsetminus}
\usepackage{xurl,hyperref}
\hypersetup{colorlinks=true,urlcolor=blue}
\setcounter{secnumdepth}{0}
\setlength{\parindent}{0pt}
\setlength{\parskip}{5pt}
\setlength{\emergencystretch}{3em}
\begin{document}
\section*{Largest common subcladogram exponents}\label{9700019}
{\small UnsolvedMath ID 9700019 · AMR-096-0019 · Probability · AMR Open Problem Lists}
\subsection{Definitions and mathematical statement}
A rooted binary cladogram on $[n]=\{1,\ldots,n\}$ is a finite rooted tree whose leaves carry these distinct labels, whose internal vertices have two unordered children, and whose internal vertices carry no labels. For $A\subseteq[n]$, its induced cladogram retains the ancestry connecting the leaves of $A$, takes their most recent common ancestor as root, and suppresses vertices with only one retained child.

Consider separately two probability laws: the uniform law on the $(2n-3)!!$ such cladograms, and the coalescent law obtained by starting with $n$ single-leaf lineages and repeatedly choosing uniformly a pair of current lineages and joining them under a new parent.

Under either law, sample $T_1,T_2$ independently and set
\[
C_n=\max\{|A|:A\subseteq[n],\ T_1|_A\cong T_2|_A\},
\]
where the isomorphism must preserve the root and each leaf label. Prove that there are constants $\gamma_a,\gamma_b<1/2$, respectively for the uniform and coalescent laws, for which
\[
\mathbb EC_n=n^{\gamma_a+o(1)},\qquad
\mathbb EC_n=n^{\gamma_b+o(1)},
\]
and identify or characterize these exponents.

The statements are equivalent to convergence of $\log\mathbb EC_n/\log n$ to the corresponding constants. The same subset of labels must be used in both trees; arbitrary relabeling is not permitted. The requested result concerns the optimum common induced cladogram, not merely the output of one particular recursive matching algorithm.
\subsection{Short English statement}
What are the growth exponents of the largest common labeled subcladogram in two independent uniform or coalescent random trees?
\subsection{Sources}
\begin{itemize}
\item[S1] ULAM.AI and the UnsolvedMath Contributors, supplied problems.json, record 9700019 (AMR-096-0019), AMR Open Problem Lists. Catalogue identity and source transcription; CC BY 4.0.
\url{https://huggingface.co/datasets/ulamai/UnsolvedMath}.
\item[S2] Aldous - Open problems, substructures.html; linked 2003 document, Example 3. Original source indexed by AMR-096-0019.
\url{https://www.stat.berkeley.edu/~aldous/Research/OP/index.html}.
\item[S3] D. Aldous, Largest common substructures in probabilistic combinatorics, original problem definitions.
\url{https://www.stat.berkeley.edu/~aldous/Research/OP/common_subs.pdf}.
\end{itemize}
\end{document}
